% RESEARCH_DRAFT_SHA256: 27ffac22addf2f3bd855e481698df0cb7707e74c1f8ff32e551fbfe39cec6d3a
% Layout and source-reference wrappers expanded for the web renderer.
\documentclass[11pt]{article}
\usepackage[margin=25mm]{geometry}
\usepackage{amsthm}
\usepackage{fontspec,unicode-math}
\setmainfont{Latin Modern Roman}
\setmathfont{Latin Modern Math}
\usepackage{amsmath,amssymb,mathtools}
\usepackage{xurl,hyperref}
\hypersetup{colorlinks=true,urlcolor=blue,linkcolor=blue}
\setcounter{secnumdepth}{0}
\setlength{\parindent}{0pt}
\setlength{\parskip}{5pt}
\setlength{\emergencystretch}{3em}
\newcommand{\sref}[2]{\hyperlink{source:#1:#2}{[S#2]}}
\newcommand{\cataloguescope}[1]{\paragraph{Source question.}#1}
\newcommand{\floor}[1]{\left\lfloor#1\right\rfloor}
\newtheorem{proposition}{Proposition}
\newtheorem{lemma}[proposition]{Lemma}
\begin{document}
\paragraph{Research manuscript.} OpenAI \texttt{gpt-6-astra}, reasoning \texttt{ultra}; 11 October 2026.
\paragraph{Preserved source body.} The complete catalogue mathematical text and sources follow with layout and source links expanded, including source qualifications and its historical preparation notes. New research begins at the section headed Research attempt.

\section{Complexity of Approximate Counting in Spin Systems}\label{statement}

\subsection{Definitions and mathematical statement}

\paragraph{Shared notation.}Standard number systems are denoted by $\mathbb N,\mathbb Z,\mathbb Q,\mathbb R,\mathbb C$; index ranges, computational models and specialized notation follow the cited problem source and the definitions below.


For a graph $G=(V,E)$, the ferromagnetic $q$-state Potts partition function is $Z_{\mathrm P}=\sum_{\sigma:V\to\{1,\ldots,q\}}\beta^{|\{uv\in E:\sigma(u)=\sigma(v)\}|}$ with $\beta>1$. The hard-core partition function is $Z_{\mathrm H}=\sum_{I\subseteq V\,\mathrm{independent}}\lambda^{|I|}$ with fugacity $\lambda>0$. Their probability distributions normalize the corresponding weights. Approximate counting seeks relative-error estimates for $Z$; approximate sampling seeks distributions close in total variation to these normalized weights, with polynomial dependence on input size and the requested accuracy.
\paragraph{Statement recorded in the supplied source.}
Determine the worst-case complexity of approximate counting and sampling for the ferromagnetic Potts model and for the hard-core model on bipartite graphs. Do these problems admit polynomial-time algorithms, or are they NP-hard? \href{https://doi.org/10.4171/owr/2019/19}{[S2]}
\paragraph{Scope qualification.}
This asks for a classification across model parameters and allowed graph families, rather than asserting one hardness result for all temperatures or fugacities. The Potts graph family and the precise hardness reduction must be specified in an answer.
\paragraph{Source status.}
The supplied catalogue labels this record open. Its source-status metadata was checked against the supplied snapshot on 2026-10-11. This import does not independently establish current open status or mathematical correctness. The catalogue literature assessment is dated 2026-08-22.
\subsection{Short English statement}
Classify the complexity of estimating partition functions and sampling their associated spin configurations.
\subsection{Sources}
\begin{itemize}
\item[S1] Ulam AI and the UnsolvedMath Contributors, UnsolvedMath: A Curated Collection of Open Mathematics Problems, catalogue entry. Catalogue curation and metadata: CC BY 4.0. The supplied source question is retained; definitions and the English summary are editorial additions.
\url{https://huggingface.co/datasets/ulamai/UnsolvedMath}.
\item[S2] Combinatorics, Probability and Computing (2020), Oberwolfach Reports, DOI 10.4171/owr/2019/19, pp. 1113--1189
\url{https://doi.org/10.4171/owr/2019/19}.
\item[S3] Combinatorics, Probability and Computing, Oberwolfach Reports 19/2019, primary report PDF
\url{https://ems.press/content/serial-article-files/46798}.
\textit{Underlying Oberwolfach Reports material retains its source-specific licence, including CC BY-SA 4.0 where applicable; see THIRD\_PARTY\_NOTICES.md.}
\end{itemize}
\paragraph{Preparation and review.}
Original question attributed to the cited source and UnsolvedMath contributors. Codex (GPT-6; exact serving revision not exposed) assisted with definitions, English summaries, source inspection and formatting. No new human mathematical authorship or independent proof verification is asserted. Catalogue identity checked against the complete supplied registry and existing repository statements; known flagged repetitions and obvious identical conjectures excluded. Full mathematical-equivalence/literature review has not been performed. Statement preparation and final review followed contribution prompts C and D; automated checks verify structure and identity, not open status.
\label{end:statement}

\section{Research attempt}
\subsection{Literature, source context and scope}
The full original catalogue body and source qualifications are preserved above. The relevant primary workshop passages were read in the cached publisher report. Historical claims and subsequent attributed results are distinguished from the mathematical work below. A restricted example, standard lemma or known theorem is not counted as new progress on the entire original target. Missing assumptions and repeated or rhetorical extracts are not silently turned into different solved problems.

\subsection{Exact forest recursions and their first cyclic obstruction}
The source asks for worst-case complexity beyond the graph classes covered by its low-temperature algorithms. We test the simplest local elimination route. Let $F$ be a forest with $n$ vertices and $c$ connected components. For the hard-core model use weight $\lambda^{|I|}$ for independent sets, with fixed $\lambda>0$. For a $q$-state ferromagnetic Potts model use weight $b^{m(\sigma)}$, where $q\ge2$, $b>1$, and $m(\sigma)$ counts monochromatic edges. This is equivalent, up to the known factor $b^{|E|}$, to the source's weight $e^{-\beta\,\#\text{bichromatic}}$ when $b=e^\beta$.
\begin{proposition}
Both partition functions and exact conditional sampling distributions on forests are computable with $O(n)$ arithmetic operations for fixed parameters. In particular,
\[
Z_{\mathrm{Potts}}(F)=q^c(b+q-1)^{n-c}.
\]
\end{proposition}
\begin{proof}
Root each tree. For hard-core configurations in the subtree at $v$, let $Z_v^0,Z_v^1$ be the total weights with $v$ respectively absent or present. Conditional on the state of $v$, distinct child subtrees have no connecting edges. Hence
\[
Z_v^0=\prod_{w\text{ child of }v}(Z_w^0+Z_w^1),\qquad
Z_v^1=\lambda\prod_{w\text{ child of }v}Z_w^0.
\]
For a leaf these are $1,\lambda$, and induction proves the formulas and the root total $Z_v^0+Z_v^1$. Multiplying component totals gives $Z(F)$. Choose the root present with probability $Z_v^1/(Z_v^0+Z_v^1)$. A child is absent if its parent is present; otherwise use the analogous subtree ratio. The same conditional factorization proves by induction that this samples exactly the normalized weights.
For Potts, with the parent color fixed, summing over a leaf color contributes $b+q-1$, independently of that color. Repeated leaf deletion leaves one freely colored root per component and $n-c$ deleted edges, proving the formula. Sample each root uniformly and each child with its parent's color with probability $b/(b+q-1)$ and with each other color with probability $1/(b+q-1)$. Multiplying these transition probabilities produces $b^{m(\sigma)}/Z(F)$. Each edge is processed a constant number of times in both models. For fixed rational parameters the intermediate numerators and denominators have $O(n)$ bits, so standard rational arithmetic also gives polynomial bit complexity. Exact random choices have their usual random-bit implementation; the operation bound concerns the weight computations and conditional probabilities.
\end{proof}
The attempted extension deletes an edge to apply tree recursions. It already changes the hard-core answer on a four-cycle: direct enumeration gives $Z_{C_4}=1+4\lambda+2\lambda^2$, while the path $P_4$ obtained by deleting one edge has $Z_{P_4}=1+4\lambda+3\lambda^2$. The newly allowed pair consists of the endpoints of the deleted edge. Thus a local independence factorization cannot be reused without tracking the dependency reintroduced by that edge.
\subsection{Verification and remaining gaps}
All independent sets and all three-color assignments on small forests were enumerated and matched the recursions; the cycle/path discrepancy was checked exactly. Forests are a tractable special class even at large activity or coupling. No algorithm for arbitrary bipartite graphs, no hardness classification for approximate counting, and no resolution of the source's general low-temperature agenda follows.

\section{References and source attribution}
Primary workshop report: \url{https://doi.org/10.4171/owr/2019/19}. Relevant passages read: 2019/19 pp. 1162--1163, low-temperature algorithms and the open worst-case counting/sampling classification.

All retrievals were checked on 11 October 2026. Original question attribution remains in the preserved Sources section. External theorems are attributed and not claimed as new discoveries or fully verified merely by citation.


\paragraph{Integration review.} Codex (GPT-6; exact serving revision not exposed) reviewed the written arguments and expanded layout and source-reference wrappers. Repository metadata and the production web renderer were checked using Python 3.13 and JavaScript with MathJax 3.2.2. This model-assisted review is not independent human proof verification; no PDF compilation or formal proof checking is asserted.

\section{Completion Estimate}
\noindent\textbf{COMPLETION: 0\%}
No substantiated new progress toward the entire original target is claimed. The worked results and failed route are delimited above; a known result, illustrative model, or auxiliary obstruction is not counted as original-target completion.

\section{Final status and attribution}
\textbf{Attempt status: unresolved.} This describes this research attempt, not a new ruling on the source catalogue's status or on external literature claims.
The mathematics and this self-audit were generated by OpenAI \texttt{gpt-6-astra}, reasoning \texttt{ultra}, on 11 October 2026. Codex (GPT-6; exact serving revision not exposed) orchestrated the batch, reviewed the written mathematical arguments, and contributed formatting/repository review. This additional model review is not independent human verification. No human mathematical authorship, independent human proof checking, or formal verification is claimed. The historical preparation paragraphs in the reproduced catalogue body concern that body only.

Contribution prompt A was applied to the complete definition and supplied source record, followed by the adversarial checks of prompt B. Local reproduction files are \texttt{batch022/checks.py}, \texttt{computations.json}, and the per-ID audit JSON. They were not public links at preparation time. Numerical checks are finite and do not prove unrestricted results. No LaTeX engine was found locally; compilation was not performed by this research agent.

\end{document}
